\section{Discussion}
\label{sec:discussion}
The shape of the dissipation matters as much as its rate, and three design rules follow for the reservoirs we
studied:
\begin{enumerate}[label=(\arabic*),leftmargin=*,itemsep=2pt,topsep=4pt]
\item let the jump operators set the direction of the clock: separate channels slow it with weaker loss, a shared
  channel with dark states slows it with stronger loss, and its phases decide whether the input reaches the dark
  states;
\item let the input follow the clock, weakening it as the memory grows so that the mixing fraction of one input times
  the memory time stays of order one;
\item as a heuristic, forget by cooling, not by randomizing: prefer dissipation that leaves a polarized state the
  input can tilt.
\end{enumerate}
A damping sweet spot found at a fixed input scaling~\cite{Sannia2024} is therefore not a property of the dissipation
alone: it moves as the input weakens, and with finite sampling the measurement budget sets it. The right shape need
not be exotic: with a uniform input, weak local loss keeps the most memory at every budget we tried.

Dark states pay when the reservoir must dissipate strongly, as in waveguides, where collective decay comes with the
coupling~\cite{Zanner2022,Sheremet2023} and the positions of the spins set the phases of the channel: spins half a
wavelength apart couple with opposite signs~\cite{Sheremet2023}. Ordinary noise limits their reach: local loss or
dephasing at a small rate adds to the leak out of the dark states, which the collective loss suppresses only as
$1/\gamma$, so a memory of $T$ steps needs collective loss about $4T^2$ times stronger than local loss
(Appendix~\ref{app:numerics}). A dark state of four transmons in a waveguide lives more than a hundred times longer
than a single qubit decaying into the guide~\cite{Zanner2022}; read as such a ratio of rates, this allows a memory of
about five steps, and the fifty steps of Fig.~\ref{fig:clock}b would need a ratio near $10^4$.

Parts of this story have classical cousins: two oscillators sharing one damper keep an undamped mode, and the
$1/\gamma$ leak is the slow mode of an overdamped oscillator. The quantum structure adds three things: the protected
set is a space of operators, 100 of them for ten dark states; the memory is written into coherences between dark
energy levels; and the phases of the jump operators and the input decide whether the input reaches the protected
states at all.

Our results hold for one class of Hamiltonians, eight reservoirs and a uniform input, whose symmetry decides which
channels write gently by themselves; with the input on a single spin the collapses persist
(Appendix~\ref{app:numerics}). We study memory and capacities, which set what a linear readout can
reproduce~\cite{Dambre2012}, not a particular task. Jumps beside a transverse drive assume a bath spectrum that is
flat over the relevant frequencies~\cite{Cattaneo2019}; an open superconducting waveguide provides one up to
$\gamma\approx64$ for couplings of order a megahertz~\cite{Sheremet2023}, and shared resonators, waveguides and
structured couplers can build such channels~\cite{Cattaneo2019,Brehm2021,Sheremet2023,CattaneoPRX2023}. Concurrent
work finds that decoherence-free modes break the echo-state property and that responses below the regularization
scale do not count~\cite{Cao2026}; our dark states leak, and Sec.~\ref{sec:price} shows the second effect. Another
finds a balance between memory and nonlinearity across reservoir protocols~\cite{Cindrak2026}, which the mixing law
makes quantitative for slow memories. Truncated replay~\cite{Cindrak2024} or weak measurements~\cite{Mujal2023}
would cut the cost of sampling; how a slow clock combines with non-Markovian memory~\cite{Sannia2026Memory} is
open.

How a quantum system leaks decides not only how fast it forgets, but what it keeps and in what form.
